% Generated draft: requires mathematical and source review.
\documentclass[11pt]{article}
\usepackage[margin=25mm]{geometry}
\usepackage{fontspec}
\setmainfont{Latin Modern Roman}
\usepackage{amsmath,amssymb,mathtools,unicode-math}
\setmathfont{Latin Modern Math}
\usepackage{enumitem,needspace,xurl,hyperref}
\setlength{\parindent}{0pt}
\setlength{\parskip}{5pt}
\setlength{\emergencystretch}{3em}
\newcommand{\N}{\mathbb{N}}
\newcommand{\Z}{\mathbb{Z}}
\newcommand{\Q}{\mathbb{Q}}
\newcommand{\R}{\mathbb{R}}
\newcommand{\C}{\mathbb{C}}
\newcommand{\F}{\mathbb{F}}
\newcommand{\A}{\mathbb{A}}
\newcommand{\PP}{\mathbb{P}}
\newcommand{\E}{\mathbb{E}}
\newcommand{\eps}{\varepsilon}
\newcommand{\dd}{\,\mathrm{d}}
\begin{document}
\section*{UnsolvedMath ID: 701238}
\textbf{Source problem code:} catalogue entry
\subsection{Definitions and mathematical statement}
For a group presentation, every listed relator is set equal to the identity. Write \([x,y]=x^{-1}y^{-1}xy\). Define
\[
\begin{aligned}
G_3=\langle h,a,P,Q\mid {}&h^3,a^2,P^2,Q^2,[P,Q],[h,P],(hQ)^2,a^{-1}PaQ\rangle,\\
G_4=\langle h,a,p,q,r\mid {}&h^3,a^2,p^2,q^2,r^2,[p,q],[p,r],p(qr)^2,\\
&h^{-1}phq,h^{-1}qhpq,(hr)^2,[a,p],a^{-1}qar\rangle.
\end{aligned}
\]
The supplied record states that both groups act arc-transitively on the infinite \(3\)-valent tree with finite vertex stabilizers. Here every vertex of the tree has degree \(3\), and arc-transitivity means transitivity on ordered pairs of adjacent vertices. It also states that each group contains
\[
G_1=\langle h,a\mid h^3,a^2\rangle\cong\operatorname{PSL}_2(\Z)
\]
as a subgroup of finite index.

Form the amalgamated free product \(G=G_3*_{G_1}G_4\), identifying the displayed copies of \(h\) and \(a\). According to the supplied record, \(G\) has a normal subgroup \(K\) of index \(8\), generated by
\[
A=h,\qquad B=aha,\qquad C=p,\qquad D=PpP,\qquad E=QpQ,\qquad F=PQpPQ,
\]
with dihedral complement \(\langle a,P,Q\rangle\); thus the asserted decomposition is \(G=K\rtimes\langle a,P,Q\rangle\), with this complement of order \(8\). The supplied presentation of \(K\), which specifies the group asked about, is
\[
\begin{aligned}
K=\langle A,B,C,D,E,F\mid {}&A^3,B^3,C^2,D^2,E^2,F^2,\\
&(AC)^3,(AD)^3,(AE)^3,(AF)^3,\\
&(BC)^3,(BD)^3,(BE)^3,(BF)^3,\\
&(ABA^{-1}C)^2,(ABA^{-1}D)^2,\\
&(A^{-1}BAE)^2,(A^{-1}BAF)^2,\\
&(BAB^{-1}C)^2,(B^{-1}ABD)^2,\\
&(BAB^{-1}E)^2,(B^{-1}ABF)^2\rangle.
\end{aligned}
\]
\textbf{Question.} Does there exist a finite group \(H\) with \(|H|>1\) and a surjective homomorphism \(K\twoheadrightarrow H\)? Equivalently, does there exist \(N\trianglelefteq K\) such that \(1<[K:N]<\infty\)? The structural assertions preceding the question are supplied source assertions and have not been independently verified here.
\subsection{Short English statement}
Does the explicitly presented group \(K\), described in the supplied record as a normal subgroup of index \(8\) in \(G_3*_{G_1}G_4\), have a nontrivial finite quotient? The registry labels the problem open; that label is supplied metadata, not an independently verified assessment.
\subsection{Sources}
\begin{itemize}
\item Local UnsolvedMath registry, permanent ID \(701238\), exact source code \texttt{catalogue entry}, title \emph{Kourovka Notebook Problem 16.46}. The supplied record gives upstream revision \texttt{d371ee917ac9647fbb06cd6fc529deec6f3dda82} and an update timestamp of \texttt{2026-10-16T14:00:00Z}.
\item M. Conder, Problem 16.46, \emph{The Kourovka Notebook (Unsolved Problems in Group Theory)}, 16th issue (2006), as attributed in the supplied record. The record also identifies the 21st issue edition (2026), \texttt{arXiv:1401.0300v48}. Its collection description names E. I. Khukhro and V. D. Mazurov as editors. These bibliographic details are supplied attribution; the cited edition has not been independently consulted here.
\end{itemize}
\end{document}
